% !PRINTED 434
\weberpage{434}
%<<CONT
and in it, since $\tau$ is to be of degree 7, the $\varepsilon_1, \varepsilon_2, \varepsilon_3$
must be seventh roots of unity which satisfy the condition
\weq{2}{\varepsilon_1 \varepsilon_2 \varepsilon_3 = 1}

We shall now seek to determine the substitution $\chi$ so
that the condition $\chi \tau = \tau^4 \chi$, or
\weq{3}{\tau \chi = \chi \tau^2}
is fulfilled. If we take $\chi$ in the form
\[
\chi = \begin{pmatrix} \alpha_1, & \alpha_2, & \alpha_3 \\ \beta_1, & \beta_2, & \beta_3 \\ \gamma_1, & \gamma_2, & \gamma_3 \end{pmatrix}
\]
then the condition (3) yields
\[
\begin{pmatrix} \alpha_1 \varepsilon_1, & \alpha_2 \varepsilon_1, & \alpha_3 \varepsilon_1 \\ \beta_1 \varepsilon_2, & \beta_2 \varepsilon_2, & \beta_3 \varepsilon_2 \\ \gamma_1 \varepsilon_3, & \gamma_2 \varepsilon_3, & \gamma_3 \varepsilon_3 \end{pmatrix}
=
\begin{pmatrix} \alpha_1 \varepsilon_1^2, & \alpha_2 \varepsilon_2^2, & \alpha_3 \varepsilon_3^2 \\ \beta_1 \varepsilon_1^2, & \beta_2 \varepsilon_2^2, & \beta_3 \varepsilon_3^2 \\ \gamma_1 \varepsilon_1^2, & \gamma_2 \varepsilon_2^2, & \gamma_3 \varepsilon_3^2 \end{pmatrix}.
\]

Now since $\alpha_1, \beta_1, \gamma_1$ cannot all three vanish, $\varepsilon_1^2$
must agree with one of the three quantities $\varepsilon_1, \varepsilon_2, \varepsilon_3$.
If we had $\varepsilon_1 = \varepsilon_1^2$, hence $\varepsilon_1 = 1$, and consequently $\varepsilon_3 = \varepsilon_2^{-1}$, then
we could have neither $\varepsilon_2$ nor $\varepsilon_3 = 1$, since otherwise $\tau$ would be the identical
substitution. We should therefore have to have $\beta_1 = \gamma_1 = \alpha_2 = \alpha_3 = 0$ and $\varepsilon_2^2 = \varepsilon_3$,
$\varepsilon_2^3 = 1$, which is impossible. There remain, therefore, by (2),
two possible assumptions:
\[
\varepsilon_1 = \varepsilon, \quad \varepsilon_2 = \varepsilon^2, \quad \varepsilon_3 = \varepsilon^4,
\]
\[
\varepsilon_1 = \varepsilon, \quad \varepsilon_2 = \varepsilon^4, \quad \varepsilon_3 = \varepsilon^2,
\]
in which $\varepsilon$ is an imaginary seventh root of unity, and these
two assumptions lead to two \gesperrt{different} groups, which are
built quite alike, and moreover are also isomorphic. We pursue the first assumption further here, and accordingly set
\weq{4}{\tau = \begin{pmatrix} \varepsilon, & 0, & 0 \\ 0, & \varepsilon^2, & 0 \\ 0, & 0, & \varepsilon^4 \end{pmatrix}, \quad \chi = \begin{pmatrix} 0, & 0, & 1 \\ 1, & 0, & 0 \\ 0, & 1, & 0 \end{pmatrix}, \quad \chi^2 = \begin{pmatrix} 0, & 1, & 0 \\ 0, & 0, & 1 \\ 1, & 0, & 0 \end{pmatrix}, \quad \chi^3 = 1.}

That we have set the non-vanishing quantities $\alpha_3, \beta_1, \gamma_2$ equal to 1
is no restriction, since we can reduce every other assumption to this one by a transformation of the whole
group.

The meaning of $\chi$ in this form is that of a cyclic
permutation of the variables. The substitutions $\tau, \chi$ together
generate a group $\tau^\varrho \chi^\lambda$ of degree 21.
